\documentclass[11pt,a4paper]{article}
\usepackage{templates/modeling-notes}
\hypersetup{pdftitle={Simplex practice -- answer key}}
\title{\textbf{Mathematical Modeling}\\[2mm]\large Simplex practice: answer key}
\author{Personal study notes}
\date{}
\begin{document}
\maketitle
Attempt the separate practice sheet before reading this document. The results
below are exact mathematical answers. Run the accompanying reference script
only after completing your own MATLAB models.

\section{Mock lab A: slack basis and two pivots}
The original model is
\[
  \max(5x_1+4x_2),\qquad
  \begin{bmatrix}3&2\\1&2\end{bmatrix}x\leq
  \begin{bmatrix}18\\10\end{bmatrix},\qquad x\geq0.
\]
For MATLAB, append $-I_2x\leq0$ to the two resource rows and minimize
$(-5,-4)x$. Thus $A$ is $4\times2$, $b$ has four entries and $f$ has two.
For standard form, with $v=(x_1,x_2,s_1,s_2)^T$, use
\[
  S=\begin{bmatrix}3&2&1&0\\1&2&0&1\end{bmatrix},\quad
  b=\begin{bmatrix}18\\10\end{bmatrix},\quad c=(5,4,0,0)^T.
\]
The slack BFS is $(0,0,18,10)^T$. Initially $z_j-c_j=(-5,-4,0,0)$.
Enter $x_1$, remove $s_1$, and obtain
\[
  \left[\begin{array}{rrrr|r}
      1&2/3&1/3&0&6\\0&4/3&-1/3&1&4
    \end{array}\right],\qquad
  z_j-c_j=(0,-2/3,5/3,0),\quad z=30.
\]
Now enter $x_2$ and remove $s_2$. The final constraint rows are
\[
  \left[\begin{array}{rrrr|r}
      1&0&1/2&-1/2&4\\0&1&-1/4&3/4&3
    \end{array}\right],\qquad
  z_j-c_j=(0,0,3/2,1/2),\quad z=32.
\]
Thus $x^*=(4,3)^T$, both slacks vanish, and profit is 32. The weighted
$z_j$ rows, including the objective value on the RHS, are respectively
$(0,0,0,0\mid0)$, $(5,10/3,5/3,0\mid30)$ and
$(5,4,3/2,1/2\mid32)$.

The feasible polygon has vertices $(0,0),(6,0),(4,3),(0,5)$; the hand path
visits its first three vertices in that order. A direct certificate is
\[
  5x_1+4x_2=\tfrac32(3x_1+2x_2)+\tfrac12(x_1+2x_2)
  \leq\tfrac32(18)+\tfrac12(10)=32.
\]
For 20 assembly hours, the same two resource constraints are tight at
$x^*=(5,5/2)^T$, with profit 35. The same nonnegative resource weights
give an upper bound of 35, proving optimality. Extra assembly capacity shifts
production toward the product with greater assembly use, while finishing remains
fully used. In MATLAB the standard model uses $S$ as equalities,
$-I_4v\leq0$ as inequalities, and objective $-c$.

\clearpage
\section{Mock lab B: two-phase initialization}
The original model is
\[
  \min C=2x_1+3x_2,\qquad x_1+x_2\geq4,\quad x_1+3x_2=6,
  \qquad x_1,x_2\geq0.
\]
For the original MATLAB call, use
\[
  f=\begin{bmatrix}2\\3\end{bmatrix},\quad
  A=\begin{bmatrix}-1&-1\\-1&0\\0&-1\end{bmatrix},\quad
  b=\begin{bmatrix}-4\\0\\0\end{bmatrix},\quad
  A_{\mathrm{eq}}=\begin{bmatrix}1&3\end{bmatrix},\quad b_{\mathrm{eq}}=6.
\]
The original standard form is
\[
  \begin{bmatrix}1&1&-1\\1&3&0\end{bmatrix}
  \begin{bmatrix}x_1\\x_2\\s_1\end{bmatrix}
  =\begin{bmatrix}4\\6\end{bmatrix},\qquad x_1,x_2,s_1\geq0.
\]
Setting $x_1=x_2=0$ would give $s_1=-4$ and would also violate the
second equality. It is not a feasible start. For Phase I, append $a_1,a_2$
and maximize $-a_1-a_2$:
\[
  \left[\begin{array}{rrrrr|r}
      1&1&-1&1&0&4\\1&3&0&0&1&6
    \end{array}\right].
\]
With $B=(a_1,a_2)$, $c_B=(-1,-1)^T$ and $z_{\mathrm I}=-10$, the initial
indicator vector is $(-2,-4,1,0,0)^T$. Enter $x_2$ and remove $a_2$:
\[
  \left[\begin{array}{rrrrr|r}
      2/3&0&-1&1&-1/3&2\\1/3&1&0&0&1/3&2
    \end{array}\right],\quad
  z_j-c_j=(-2/3,0,1,0,4/3)^T,\quad z_{\mathrm I}=-2.
\]
Enter $x_1$ and remove $a_1$. Apply $R_1\leftarrow(3/2)R_1$, then
$R_2\leftarrow R_2-(1/3)R_1$ using the new row 1:
\[
  \left[\begin{array}{rrrrr|r}
      1&0&-3/2&3/2&-1/2&3\\0&1&1/2&-1/2&1/2&1
    \end{array}\right].
\]
Both artificials are nonbasic and zero. The Phase I indicators are
$(0,0,0,1,1)^T$, and the optimal artificial sum is zero. Hence the
original system is feasible. Delete the last two columns and recompute using
the original maximization coefficients $(-2,-3,0)^T$:
\[
  z=-C=-9-\tfrac32s_1,\qquad z_j-c_j=(0,0,3/2)^T.
\]
No Phase II pivot is needed. The unique minimum is $C^*=9$ at $(3,1)^T$.
As an independent check, eliminate $x_1=6-3x_2$. Nonnegativity and the
minimum total mass give $0\leq x_2\leq1$, and
$C=12-3x_2$ is smallest at $x_2=1$.

For MATLAB's original standard form, minimize $(2,3,0)^Tv$, use the
three-column matrix above as equalities, and append $-I_3v\leq0$.
Artificial variables are a device for our hand initialization; the solver
handles its own initialization. The reference script separately solves the
explicit Phase I problem to check that its artificial sum is zero, without
assuming that the solver follows these particular hand pivots.

\clearpage
\section{Mock lab C: signs, shifts and recovery}
The original MATLAB inequalities and minimized objective are
\[
  A=\begin{bmatrix}1&1\\-1&1\\0&-1\\0&1\end{bmatrix},\qquad
  b=\begin{bmatrix}3\\1\\1\\2\end{bmatrix},\qquad
  f=\begin{bmatrix}-3\\-2\end{bmatrix}.
\]
The third row is $v\geq-1$, rather than $v\geq0$. No lower bound is imposed
on $u$. Recover the full score as $-\texttt{fval}+5$.

Let $u=u^+-u^-$ and $v=w-1$. The three upper inequalities become
\[
  u^+-u^-+w\leq4,\qquad -u^++u^-+w\leq2,\qquad w\leq3.
\]
The objective is $F=3u^+-3u^-+2w+3$. With standard-form vector
$y=(u^+,u^-,w,s_1,s_2,s_3)^T\geq0$,
\[
  S=\begin{bmatrix}
    1&-1&1&1&0&0\\-1&1&1&0&1&0\\0&0&1&0&0&1
  \end{bmatrix},\quad
  b=\begin{bmatrix}4\\2\\3\end{bmatrix},\quad c=(3,-3,2,0,0,0)^T.
\]
The slack start $y=(0,0,0,4,2,3)^T$ represents $(u,v)=(0,-1)$,
with full score 3. Initially $z_j-c_j=(-3,3,-2,0,0,0)^T$ for the
variable part of the objective. Enter $u^+$ and remove $s_1$. Adding row 1
to row 2 gives
\[
  \left[\begin{array}{rrrrrr|r}
      1&-1&1&1&0&0&4\\0&0&2&1&1&0&6\\0&0&1&0&0&1&3
    \end{array}\right],\qquad z_j-c_j=(0,0,1,3,0,0)^T.
\]
The variable objective is 12, so the full score is $12+3=15$. One
representative optimum is $y=(4,0,0,0,6,3)^T$, recovering
$(u,v)=(4,-1)$. A direct upper-bound certificate is
\[
  F=3(u+v)-v+5\leq3(3)+1+5=15.
\]
Equality forces $u+v=3$ and $v=-1$, so the original optimum is unique.
The split representation is not unique: replacing $(u^+,u^-)$ by
$(4+t,t)$ for any $t\geq0$ leaves $u$ and the objective unchanged.
For example $(6,2)$ represents the same signed adjustment as $(4,0)$.
The zero indicator for $u^-$ reflects this redundancy, not a second original
decision. Compare recovered variables, not an arbitrarily chosen split.

For MATLAB's standard form, minimize $-c^Ty$ subject to $Sy=b$ and
$-I_6y\leq0$. Recover the full score as $-\texttt{fval}+3$.
The different constants in the original and transformed calls are essential.

\clearpage
\section{Mock lab D: three different outcomes}
\subsection{D1: an optimal segment}
In standard form append $I_3$ to the original three-row matrix and use
$(x_1,x_2,s_1,s_2,s_3)\geq0$. Its slack start is $(0,0,4,3,3)^T$.
The first restriction gives the bound $x_1+x_2\leq4$. All optimizers are
\[
  (x_1,x_2)=(t,4-t),\qquad1\leq t\leq3,
\]
with objective value 4. In particular $(1,3)$ and $(3,1)$ are distinct
optimal points. Check the sum, bounds and objective of the returned vector;
do not require one endpoint. In the original MATLAB call use resource rows
$(1,1),(1,0),(0,1)$ with RHS $(4,3,3)^T$, append $-I_2x\leq0$, and
minimize $(-1,-1)x$.

\subsection{D2: infeasible, with a Phase I certificate}
The original requirements imply both $x_1+x_2\geq5$ and $x_1+x_2\leq4$,
an impossibility. For the augmented equations in the question,
the second gives $x_1+x_2=4-t/2$. Substitution into the first yields
\[
  a=1+s+t/2\geq1.
\]
The point $(x_1,x_2,s,t,a)=(4,0,0,0,1)$ attains $a=1$, so the Phase I
minimum is strictly positive. There is no original feasible point and hence
no original optimal value or optimizer. MATLAB input uses the rows
$(-1,-1)$ and $(2,2)$ with RHS $(-5,8)^T$, plus nonnegativity rows,
and objective $(1,2)x$. An infeasible status is the expected outcome.

\subsection{D3: feasible and unbounded above}
Standard form is $x_1-x_2+s=1$, with $x_1,x_2,s\geq0$. From the slack
start, enter $x_1$ and remove $s$. The point becomes $(x_1,x_2,s)=(1,0,0)$.
The objective equation is
\[
  z=2+3x_2-2s.
\]
Thus the indicator for $x_2$ is $-3$, but its transformed column is $(-1)$:
there is no positive entry and no leaving row. Set $x_2=\theta$, $s=0$,
so $x_1=1+\theta$. In original coordinates,
\[
  x(\theta)=(1+\theta,\theta)^T=(1,0)^T+\theta(1,1)^T,\qquad\theta\geq0.
\]
Its constraint left side stays at 1 and $z(\theta)=2+3\theta\to+\infty$.
The original MATLAB model minimizes $(-2,-1)x$ subject to $(1,-1)x\leq1$
and $-I_2x\leq0$. An unbounded status refers to minimization of the negative
profit; it confirms there is no finite maximum of the original profit.

\section{Reference MATLAB checks}
\href{run:matlab/labs/lab_03_practice_answers.m}{\texttt{lab\_03\_practice\_answers.m}}
contains complete solver matrices, feasibility/objective assertions and the
certificates above. Run from the repository root:
\begin{lstlisting}[language=Matlab]
run('matlab/labs/lab_03_practice_answers.m')
\end{lstlisting}
Constraint and optimality tolerances are $10^{-9}$; comparisons use absolute
tolerance $10^{-7}$. For finite optima, check a positive exit flag before using
the solution. For D2 and D3, inspect the solver message and compare it with the
independent infeasibility or ray certificate. A numerical failure is not itself
a mathematical infeasibility proof.

\section{Sources and editorial remarks}
These are independent solutions to invented practice tasks. The confirmed
broad Lab 3 topics are simplex and standard form. The questions, certificates, time budgets and answers are
not official assessment material. The numerical hand results were checked
with exact rational calculations independently of MATLAB.

The solver interface follows the official
\href{https://www.mathworks.com/help/optim/ug/linprog.html}{\texttt{linprog} reference}.
It was checked on 25 October 2026. The demo and reference scripts require Optimization
Toolbox for their solver calls. Attempted execution through the installed MATLAB
launcher ended with code 1 and no output; actual MATLAB execution and plot
rendering remain unverified. See \path{docs/matlab-verification.md}.
\end{document}
